\QHeading{第4题　脱式计算。}
\par\noindent


\CalcBlock[简便解法（补充）]{1}{\begin{aligned}[t]&\phantom{={}}\color{black}(540\div 27+39)\times 17\\&\color{answerred}=(20+39)\times 17\\&\color{answerred}=59\times 17\\&\color{answerred}=(60-1)\times 17\\&\color{answerred}=60\times 17-1\times 17\\&\color{answerred}=1020-17\\&\color{answerred}=1003\end{aligned}}
\par
\MethodBlock{按运算顺序计算（基础练习）}{本单元还要练习运算顺序；与上面的简便解法对照。}{\begin{aligned}[t]&\phantom{={}}(540\div27+39)\times17\\&=(20+39)\times17\\&=59\times17\\&=1003\end{aligned}}

\MethodBlock{方法二：乘法分配律}{分别求出两部分的积，再合起来。}{\begin{aligned}[t]&\phantom{={}}(540\div 27+39)\times 17\\
&=540\div 27\times 17+39\times 17\\
&=20\times 17+663\\
&=340+663\\
&=1003\end{aligned}}


\CalcBlock{2}{\begin{aligned}[t]&\phantom{={}}\color{black}475\div[(26+99)\div25]\\&\color{answerred}=475\div[125\div25]\\&\color{answerred}=475\div5\\&\color{answerred}=95\end{aligned}}
\CalcBlock{3}{\begin{aligned}[t]&\phantom{={}}\color{black}75-(416+536)\div34\\&\color{answerred}=75-952\div34\\&\color{answerred}=75-28\\&\color{answerred}=47\end{aligned}}
\CalcBlock[简便解法（补充）]{4}{\begin{aligned}[t]&\phantom{={}}\color{black}(320-144\div 6)\times 15\\&\color{answerred}=(320-24)\times 15\\&\color{answerred}=296\times 15\\&\color{answerred}=(300-4)\times 15\\&\color{answerred}=300\times 15-4\times 15\\&\color{answerred}=4500-60\\&\color{answerred}=4440\end{aligned}}
\MethodBlock{按运算顺序计算（基础练习）}{本单元还要练习运算顺序；与上面的简便解法对照。}{\begin{aligned}[t]&\phantom{={}}(320-144\div6)\times15\\&=(320-24)\times15\\&=296\times15\\&=4440\end{aligned}}

\MethodBlock{方法二：乘法分配律}{分别求出两部分的积，再合起来。}{\begin{aligned}[t]&\phantom{={}}(320-144\div 6)\times 15\\
&=320\times 15-144\div 6\times 15\\
&=4800-24\times 15\\
&=4800-360\\
&=4440\end{aligned}}

\CalcBlock{5}{\begin{aligned}[t]&\phantom{={}}\color{black}28\times[592\div(71-34)]\\&\color{answerred}=28\times[592\div37]\\&\color{answerred}=28\times16\\&\color{answerred}=448\end{aligned}}
\CalcBlock[方法一]{6}{\begin{aligned}[t]&\phantom{={}}\color{black}800\div[25\times(22-18)]\\&\color{answerred}=800\div[25\times4]\\&\color{answerred}=800\div100\\&\color{answerred}=8\end{aligned}}
\MethodBlock{方法二：除以积，改成连续除}{连续除以两个数，等于除以这两个数的积。}{\begin{aligned}[t]&\phantom{={}}800\div [25\times (22-18)]\\
&=800\div 25\div (22-18)\\
&=800\div 25\div 4\\
&=32\div 4\\
&=8\end{aligned}}

\par\noindent

\QHeading{第5题　根据运算顺序，把下列每组算式分别合并成综合算式。}

\Problem{（1）\quad \(36-8=28\)\quad \(20+30=50\)\quad \(28\times50=1400\)}
\par

\Answer{综合算式：$(36-8)\times(20+30)$}
\Answer{合并方法：把28替换为\(36-8\)，把50替换为\(20+30\)；差与和分别加上小括号，作为乘法的两个因数。}

\Problem{（2）\quad \(42-30=12\)\quad \(12\times5=60\)\quad \(480\div60=8\)}
\Answer{综合算式：$480\div[(42-30)\times5]$}
\Answer{合并方法：把60替换为\(12\times5\)，再把12替换为\(42-30\)；减法加小括号，整个除数加中括号。}

\par\noindent

\QHeading{第6题　解决问题。}

\Problem{①学校组织春游，租了6辆大车，每车坐了40名学生；又租了4辆小车，每车坐了25名学生。参加春游的学生一共有多少名？}
\par


\MethodTitle{方法一}
\Answer{\Formula{\begin{aligned}[t]&\phantom{={}}40\times6+25\times4\\&=240+100\\&=340\text{（名）}\end{aligned}}\par 答：参加春游的学生一共有340名。}

\MethodTitle{方法二：先把所有车按小车人数算，再补大车的差额}
\Answer{10辆车先按每车25人计，6辆大车每车再补15人。}
\Answer{\Formula{\begin{aligned}[t]&\phantom{={}}25\times (6+4)+(40-25)\times 6\\
&=25\times 10+15\times 6\\
&=250+90\\
&=340\text{（名）}\end{aligned}}}
\Answer{答：参加春游的学生一共有340名。}
\Problem{②书店运来一批图书，其中故事书有240本，科技书的数量是故事书的一半，漫画书比科技书和故事书的总数少80本。漫画书有多少本？}

\MethodTitle{方法一}
\Answer{\Formula{\begin{aligned}[t]&\phantom{={}}240+240\div2-80\\&=240+120-80\\&=360-80\\&=280\text{（本）}\end{aligned}}\par 答：漫画书有280本。}

\MethodTitle{方法二：把故事书看成两份，科技书看成一份}
\Answer{科技书120本，故事书是这样的2份，合起来是3份，再减80本。}
\Answer{\Formula{\begin{aligned}[t]&\phantom{={}}240\div 2\times (2+1)-80\\
&=240\div 2\times 3-80\\
&=120\times 3-80\\
&=360-80\\
&=280\text{（本）}\end{aligned}}}
\Answer{答：漫画书有280本。}

